% Inbox for the lecture in progress. Type today's raw notes here: main.tex \inputs
% this file after every chapter, so the Overleaf preview builds as you go without the
% material having to belong to a chapter yet, and /integrate empties it again once
% the material has been spread across whichever chapters it belongs in.

It turns out that the martingale we constructed for the chromatic number problem is a specific instance of a more general martingale, known as the Doob Martingale.

\begin{boxdefinition}[Doob Martingale]
    For random variables $X, Y_1, Y_2, \ldots$ on some probability space, define
    \begin{align*}
        X_0 := X
        \qquad
        \forall i \geq 1, \, X_i := \E\of{X \mid Y_1, \ldots, Y_i}
    \end{align*}
    It is possible to show that $X_0, X_1, X_2, \ldots$ is a martingale, called the \textbf{Doob Martingale}.
\end{boxdefinition}

What makes this worth having is that its increments are controlled as soon as $X$ depends on no one of the $Y_i$ too heavily.

\begin{boxdefinition}[Coordinatewise Lipschitz]
    Given metric spaces $\parenth{X_i, d_i}_{i=1}^{n}$ and $(Y, u)$, we say a function $f : X_1 \times \cdots \times X_n  \to Y$ is \textbf{$k$-coordinatewise-Lipschitz} if $u\of{f\of{\mathbf{a}}, f\of{\mathbf{b}}} \leq k$ if $\mathbf{a}$ and $\mathbf{b}$ differ in exactly one coordinate.
\end{boxdefinition}

Such an $f$, fed through the Doob martingale of its arguments, concentrates as sharply as a sum would.

\begin{boxproposition}[McDiarmid's Inequality]
    Suppose $Y_1, \ldots, Y_n$ are independent random variables on domains $\mathcal{Y}_1, \ldots, \mathcal{Y}_n$. Suppose $f : \mathcal{Y}_1 \times \cdots \times \mathcal{Y}_n \to \R$ is $1$-coordinatewise-Lipschitz. Denoting $X = f\of{Y_1, \ldots, Y_n}$, we have that
    \begin{align*}
        \Pr\of{X \geq \E\of{X} + a} \leq e^{-\frac{a^2}{2n}}
    \end{align*}
    for all $a \in \R$.
\end{boxproposition}
\begin{proof}
    Consider the Doob Martingale $X_i = \E\of{X \mid Y_1, \ldots, Y_i}$. Suppose we've fixed $y_1, \ldots, y_{i - 1}$. Define
    \begin{align*}
        g_i(y) = \E\of{f\of{y_1, \ldots, y_{i-1}, y, Y_{i+1}, \ldots, Y_n}}
    \end{align*}
    \sorry % [CLAUDE] there is no need to fill this \sorry for now: it is a placeholder. At some point, Wes will post notes about Azuma and McDiarmid, which I will give you to integrate into these notes (WITH ATTRIBUTION).
\end{proof}

One application of McDiarmid's Lemma is something called the Smiley Face Lemma, from a paper of Alan Frieze and Wes~\cite{FriezePegdenSubadditive}.

For a finite set $S$, any $\eps > 0$, and any $D \ssq \R$, an $\eps$-$D$ copy of $S$ in $\R^2$ is some embedding of $S$ into $D$ such that distinct points can be separated by $\eps$-balls.

Scattering $n$ points at random over a square of area $n$ leaves linearly many such copies.

\begin{boxlemma}[Smiley Face Lemma, \cite{FriezePegdenSubadditive}]
    For any fixed $S, \eps, D$ as above, there is some $\delta > 0$ such that $n$ random points ina  $\sqrt{n} \times \sqrt{n}$ square have $\delta n$ many $\eps$-$D$ copies of $S$ (with high probability).
\end{boxlemma}

\section{The \protect\Lovasz Local Lemma}

Suppose $B_1, \ldots, B_n$ are independent events with $\Pr\of{B_i} \leq 1$. It's pretty easy to see that $\Pr\of{\bigcap B_i^C} > 0$. Essentially, we think of the $B_i$ as \textit{bad events}: events we \textit{don't} want to see happen.

We want to extend this. In what way exactly? Here's a motivating question.

In a $k$-uniform hypergraph where each edge intersects at most $D$ others, what condition on $k$ makes the hypergraph $2$-colourable?

What we'd want to do is randomly colour vertices: each edge of $f$ is monochromatic with probability $\Pr\of{B_f} = \frac{1}{2^{k-1}}$. We would then like to apply some version of the fact that $\Pr\of{\bigcap B_i^C} > 0$.

Two edges sharing a vertex are plainly not independent, so what we ask instead is that each bad event depend on only a few of the others---and a graph is the natural place to record which.

\begin{boxdefinition}[Dependency Graph]
    Given a collection $\mathcal{A}$ of events, a \textbf{dependency graph on $\mathcal{A}$} is a graph with vertices $\setst{v_A}{A \in \mathcal{A}}$ such that $v_A \not\sim v_{B_1}, \ldots, v_{B_n}$ implies that $A$ is independent of $B_1, \ldots, B_n$.
\end{boxdefinition}

Another way of phrasing this is that adjacency represents dependence (ie, lack of independence).

\begin{boxlnotation}[Dependence]
    Borrowing from our use of $\sim$ for adjacency, we will denote that events $A$ and $B$ are dependent (ie, not independent) by $A \sim B$.
\end{boxlnotation}

The extension we were after weighs each event against the ones it depends on.

\begin{boxtheorem}[\Lovasz Local Lemma]
    Given a collection $\mathcal{A}$ of events and real numbers $0 < x_A \leq 1$ such that for all $A \in \mathcal{A}$,
    \begin{align*}
        \Pr\of{A} \leq x_A \prod_{B \sim A} \parenth{1 - x_B}
    \end{align*}
    we can conclude that
    \begin{align*}
        \Pr\of{\bigcap_{A \in \mathcal{A}} A^C} > 0
    \end{align*}
\end{boxtheorem}
\begin{proof}
    \sorry % [CLAUDE] we will do this next time, leave the local lemma and this theorem in this file for now
\end{proof}
